\documentclass[11pt]{article}

\usepackage[margin=1in]{geometry}
\usepackage{booktabs}
\usepackage{graphicx}
\usepackage{hyperref}
\usepackage{amsmath}
\usepackage{enumitem}
\usepackage{xcolor}
\usepackage{titlesec}
\usepackage{fancyhdr}
\usepackage{multirow}
\usepackage{caption}

\definecolor{yellow_conn}{HTML}{F9DF6D}
\definecolor{green_conn}{HTML}{A0C35A}
\definecolor{blue_conn}{HTML}{B0C4EF}
\definecolor{purple_conn}{HTML}{BA81C5}

\hypersetup{colorlinks=true, linkcolor=blue!60!black, urlcolor=blue!60!black, citecolor=blue!60!black}

\pagestyle{fancy}
\fancyhf{}
\rhead{STA 561D --- Duke University}
\lhead{Infinite Connections}
\rfoot{\thepage}

\titleformat{\section}{\large\bfseries}{\thesection}{1em}{}
\titleformat{\subsection}{\normalsize\bfseries}{\thesubsection}{1em}{}

\title{\textbf{Infinite Connections}\\[4pt]
\large AI-Generated NYT-Style Connections Puzzles at Scale}
\author{[Team Member 1] \and [Team Member 2] \and [Team Member 3]}
\date{Spring 2026}

\begin{document}
\maketitle

% =========================================================================
\section{Project Overview}
% =========================================================================

We built an end-to-end system that generates, validates, and serves
NYT \emph{Connections}-style word puzzles. The system combines large language
model creativity (Claude, Anthropic API) with sentence-transformer embeddings
(MPNET) to produce puzzles that are diverse, unambiguous, and difficulty-calibrated.

The codebase is fully functional in \texttt{DRY\_RUN} mode with no API keys
required. All 31 tests pass, both notebooks execute without errors, and the
web app serves playable puzzles from mock data.

% =========================================================================
\section{What We Built}
% =========================================================================

\subsection{Puzzle Generation Pipeline (\texttt{src/generator/})}

Each puzzle is assembled through six stages:

\begin{enumerate}[nosep]
    \item \textbf{Group creation} --- Claude proposes a category name and 8
          candidate words. A \emph{story injection} prompt (4 random seed words
          from the NYT word bank) prevents repetitive output.
    \item \textbf{Embedding selection} --- MPNET computes pairwise cosine
          similarity across all $\binom{8}{4} = 70$ subsets and selects the
          most internally cohesive 4 words.
    \item \textbf{Iteration} --- Steps 1--2 repeat 4 times with prior groups
          passed as context.
    \item \textbf{Editor pass} --- A second LLM call reviews and rewrites
          inaccurate category names.
    \item \textbf{Difficulty assignment} --- Groups are sorted by similarity
          and mapped to colors (yellow $\to$ purple).
    \item \textbf{Deduplication} --- The 16-word set is checked against 554
          known NYT puzzles; $>$6 overlapping words triggers a flag.
\end{enumerate}

We also implemented the \textbf{false-group pipeline}: a decoy category is
generated first, then each of its 4 words is reinterpreted via an alternate
meaning to seed 4 real groups. The decoy becomes a natural trap for solvers.

\subsection{Multi-Solver Validation (\texttt{src/solvers/})}

Three independent solvers verify that each puzzle has a unique solution:

\begin{itemize}[nosep]
    \item \textbf{Embedding solver} --- greedy selection from all
          $\binom{16}{4} = 1{,}820$ candidate groups by cosine similarity.
    \item \textbf{Clustering solver} --- beam search ($w=10$) maximizing
          Group Similarity $-$ Penalty:
          \[
            G = 0.4 \cdot I + 0.3 \cdot s_{\min} + 0.3 \cdot V,
            \qquad
            P = \frac{1}{|R|}\sum_{r \in R} \cos(\boldsymbol{\mu}_C, \mathbf{r})
          \]
    \item \textbf{LLM solver} --- Claude with chain-of-thought prompting
          (used as a tiebreaker).
\end{itemize}

The \textbf{Roundtable Validator} accepts a puzzle only when at least two
solvers converge to the same unique solution.

\subsection{Deliverables}

\begin{itemize}[nosep]
    \item \textbf{Jupyter notebooks} ---
          \texttt{infinite\_connections\_demo.ipynb} (pipeline demo) and
          \texttt{nyt\_data\_exploration.ipynb} (dataset analysis).
    \item \textbf{Web app} --- Flask + vanilla JS, playable at
          \texttt{localhost:5000} with tile selection, color reveals,
          ``one away'' feedback, and 4-mistake limit.
    \item \textbf{Docs} --- executive summary, FAQ, and technical appendix
          in \texttt{docs/}.
\end{itemize}

% =========================================================================
\section{NYT Dataset Analysis}
% =========================================================================

We analyzed the ground truth dataset of \textbf{554 NYT puzzles} spanning
June 2023 to December 2024.

\begin{table}[h]
\centering
\caption{NYT dataset summary statistics.}
\begin{tabular}{@{}lr@{}}
\toprule
Metric & Value \\
\midrule
Total puzzles               & 554 \\
Unique words                & 4,918 \\
Unique categories           & 2,054 \\
Multi-word entries           & 112 \\
Puzzles with difficulty      & 227 / 554 \\
Difficulty mean $\pm$ std   & $2.94 \pm 0.64$ \\
Difficulty range             & 1.3 -- 4.8 \\
\bottomrule
\end{tabular}
\end{table}

\paragraph{Category type distribution.}
We classified all 2,216 category slots with heuristic pattern matching:
Common property (46\%), Synonyms/Slang (38\%), Fill-in-the-blank (9\%),
Pop culture (6\%), and Wordplay (1\%).

\paragraph{Difficulty over time.}
Pearson correlation between puzzle order and difficulty is $r = -0.20$:
a slight trend toward \emph{easier} puzzles over time (not harder).

\paragraph{Most versatile words.}
BALL appears in 15 distinct categories, followed by RING (13), FLY (12),
LEAD (12), HEART (12), and WING (12). These high-ambiguity words are exactly
the type NYT favors for deception.

% =========================================================================
\section{Embedding Analysis}
% =========================================================================

We computed MPNET average pairwise cosine similarity for all 2,216 groups.
Since the dataset does not label colors per group, we assigned them by ranking
groups within each puzzle from highest similarity (yellow) to lowest (purple).

\begin{table}[h]
\centering
\caption{Cosine similarity by difficulty color: our measurements vs.\ paper thresholds.}
\begin{tabular}{@{}lccccc@{}}
\toprule
 & \multicolumn{2}{c}{Ranked (ours)} & & \multicolumn{2}{c}{Paper reference} \\
\cmidrule{2-3} \cmidrule{5-6}
Color & Mean & Std & & Mean & Variance \\
\midrule
\colorbox{yellow_conn}{Yellow} & 0.389 & 0.089 & & 0.40 & 0.0285 \\
\colorbox{green_conn}{Green}   & 0.315 & 0.060 & & 0.35 & 0.0214 \\
\colorbox{blue_conn}{Blue}     & 0.268 & 0.050 & & 0.29 & 0.0123 \\
\colorbox{purple_conn}{Purple} & 0.225 & 0.049 & & 0.27 & 0.0108 \\
\bottomrule
\end{tabular}
\end{table}

The rank ordering is correct and values are close to the paper thresholds.
NYT's color assignment agrees with pure similarity ranking only 43.9\% of the
time, confirming that difficulty has a knowledge/wordplay component that
embeddings cannot capture alone.

% =========================================================================
\section{Solver Benchmark Results}
% =========================================================================

We benchmarked the embedding and clustering solvers on all 554 NYT puzzles.

\begin{table}[h]
\centering
\caption{Solver performance on 554 NYT puzzles.}
\begin{tabular}{@{}lcc@{}}
\toprule
Metric & Embedding & Clustering \\
\midrule
Full puzzle solve rate & 2.0\% (11/554) & 3.4\% (19/554) \\
Avg groups correct     & 0.57 / 4        & 0.71 / 4 \\
\midrule
Yellow accuracy        & 24.5\%          & 28.0\% \\
Green accuracy         & 17.7\%          & 22.6\% \\
Blue accuracy          & 9.8\%           & 13.0\% \\
Purple accuracy        & 5.1\%           & 7.6\% \\
\midrule
Runtime (554 puzzles)  & 10.6\,s         & 525\,s \\
Per puzzle             & 0.019\,s        & 0.948\,s \\
\bottomrule
\end{tabular}
\end{table}

\paragraph{Key observations.}
\begin{itemize}[nosep]
    \item Both solvers follow the expected difficulty gradient: yellow (easiest)
          $>$ green $>$ blue $>$ purple (hardest).
    \item The clustering solver outperforms embedding across all colors, but at
          ${\sim}50\times$ the compute cost.
    \item Both solvers agreed on a full solve for only 7 puzzles (1.3\%),
          demonstrating that the two approaches find complementary solutions.
    \item Our embedding solver's 2.0\% full-solve rate is in the range reported
          in the literature ($\sim$11.6\% for a related MPNET approach).
          Differences likely stem from greedy vs.\ exhaustive search and
          multi-word entries in the NYT data.
\end{itemize}

% =========================================================================
\section{Codebase Status}
% =========================================================================

\begin{table}[h]
\centering
\caption{Implementation status.}
\begin{tabular}{@{}llc@{}}
\toprule
Component & Location & Status \\
\midrule
Configuration \& LLM client & \texttt{src/config.py}, \texttt{src/llm\_client.py} & Complete \\
Generator pipeline          & \texttt{src/generator/} (6 modules)               & Complete \\
Embedding solver            & \texttt{src/solvers/embedding\_solver.py}          & Complete \\
Clustering solver           & \texttt{src/solvers/clustering\_solver.py}         & Complete \\
LLM solver                  & \texttt{src/solvers/llm\_solver.py}               & Complete \\
Roundtable validator        & \texttt{src/solvers/roundtable.py}                & Complete \\
Evaluation metrics          & \texttt{src/evaluation/} (2 modules)              & Complete \\
Web app                     & \texttt{src/webapp/}                              & Complete \\
Pipeline demo notebook      & \texttt{notebooks/infinite\_connections\_demo.ipynb} & Complete \\
Data exploration notebook   & \texttt{notebooks/nyt\_data\_exploration.ipynb}   & Complete \\
Test suite                  & \texttt{tests/test\_pipeline.py} (31 tests)       & All passing \\
Solver benchmark            & \texttt{scripts/benchmark\_solvers.py}            & Complete \\
Mock data (5 puzzles)       & \texttt{data/mock/mock\_puzzles.json}             & Complete \\
Documentation               & \texttt{docs/}, \texttt{README.md}               & Complete \\
\bottomrule
\end{tabular}
\end{table}

% =========================================================================
\section{Next Steps}
% =========================================================================

\subsection{Phase 2: Scale and Evaluate}

\begin{enumerate}[nosep]
    \item \textbf{Live puzzle generation.}
          Run the pipeline with \texttt{DRY\_RUN=false} using the Anthropic API.
          Start with 100 candidates ($\sim$\$4 with Claude Sonnet), validate with
          the Roundtable, and inspect quality before scaling.

    \item \textbf{Batch generation at scale.}
          Generate 1,000--10,000 candidate puzzles. Add async/parallel API calls
          for throughput. Target $\sim$40\% validation pass rate based on the
          literature.

    \item \textbf{Solver tuning.}
          The clustering solver's beam width (currently 10) and the Group
          Similarity weights ($0.4, 0.3, 0.3$) can be tuned against the NYT
          ground truth to improve accuracy. Consider a hybrid solver that
          combines embedding scores with LLM reasoning.

    \item \textbf{Improve the embedding solver.}
          Explore exhaustive partition search (feasible for 16 words) instead of
          greedy selection to close the gap with the literature's 11.6\% baseline.
\end{enumerate}

\subsection{Phase 3: Polish Deliverables}

\begin{enumerate}[nosep, start=5]
    \item \textbf{Human evaluation.}
          Run a small user study (10--20 participants) comparing generated puzzles
          to NYT originals. Measure solve rate, enjoyment, and perceived
          difficulty.

    \item \textbf{Web app enhancements.}
          Add puzzle difficulty selector, timer, share-results feature, and
          leaderboard. Serve from the validated generated dataset instead of
          mock data.

    \item \textbf{Notebook polish.}
          Replace dry-run outputs with real generated puzzle examples. Add
          visualizations comparing generated vs.\ NYT similarity distributions
          side by side.

    \item \textbf{Written deliverables.}
          Finalize the 2-page executive summary, 2--5 page FAQ, and technical
          appendix. Incorporate actual generation results and benchmark numbers.
\end{enumerate}

\subsection{Stretch Goals}

\begin{enumerate}[nosep, start=9]
    \item \textbf{Adaptive difficulty.}
          Tune the false-group overlap to control solve rates: more overlap
          between the decoy and real groups $\to$ harder puzzles.

    \item \textbf{Domain-specific embeddings.}
          Fine-tune MPNET on Connections word pairs to better capture the
          non-semantic relationships (wordplay, fill-in-the-blank) that
          currently elude the embedding solvers.

    \item \textbf{LLM solver benchmark.}
          Run Claude Sonnet as a solver on all 554 NYT puzzles to establish
          the third baseline and measure the accuracy gap between
          embedding-based and reasoning-based approaches.
\end{enumerate}

% =========================================================================
\section*{References}
% =========================================================================

\begin{enumerate}[nosep, label={[\arabic*]}]
    \item A.\ Conneau et al., ``Making New Connections,'' arXiv:2407.11240, 2024.
    \item M.\ Todd et al., ``Missed Connections,'' arXiv:2404.11730, 2024.
    \item R.\ Li et al., ``Deceptively Simple,'' arXiv:2412.01621, 2024.
    \item J.\ Patel et al., ``Connecting the Dots,'' arXiv:2406.11012, 2024.
\end{enumerate}

\end{document}
